%======================================================================
% CAPÍTULO 2: MARCO TEÓRICO
%======================================================================
\chapter{Marco Te\'orico}

\section{Estado del Arte}

El middleware empresarial ha evolucionado significativamente desde sus or\'{\i}genes hasta convertirse en plataformas complejas de integraci\'on. El mercado ofrece soluciones como IBM Integration Bus, MuleSoft, Apache Camel y Spring Integration.

Sin embargo, la mayor\'{\i}a fueron dise\~nadas para mercados desarrollados. Esta brecha crea oportunidad para soluciones localmente desarrolladas.

\begin{definicion}
\item[Estado del Arte] Conjunto de tecnolog\'{\i}as, m\'etodos y pr\'acticas m\'as avanzadas para resolver un problema espec\'{\i}fico en un momento determinado.
\end{definicion}

\section{Fundamentos de Middleware}

El middleware es un software que act\'ua como puente entre aplicaciones y el sistema operativo, facilitando la comunicaci\'on e integraci\'on. Un Enterprise Service Bus (ESB) representa la evoluci\'on hacia una arquitectura basada en servicios.

Sandra Server implementa estos conceptos utilizando Go, proporcionando un ESB de alto rendimiento capaz de manejar miles de conexiones simult\'aneas.

\section{Arquitectura de Microservicios}

La arquitectura de microservicios construye aplicaciones como servicios peque\~nos, aut\'onomos y dbilmente acoplados. Cada microservicio es responsable de una funcionalidad espec\'{\i}fica y puede desplegarse independientemente.

El ecosistema Sandra implementa esta arquitectura: Server (middleware central), Consola (interfaz), SDC (escritorio), y Sentinel (procesamiento especializado).

\section{Seguridad de Confianza Cero}

El modelo Zero Trust elimina la confianza impl\'{\i}cita, verificando cada solicitud independientemente de su origen. Los principios incluyen: verificaci\'on expl\'{\i}cita de usuarios y dispositivos, acceso con m\'{\i}nimo privilegio, y verificaci\'on continua de seguridad.

Sandra Ecosystem implementa estos principios en todas las capas: cada solicitud de API es autenticada y autorizada, las comunicaciones usan TLS m\'utuo, y las sesiones tienen duraci\'on limitada.

\section{Tecnolog\'{\i}as Emergentes}

El ecosistema utiliza tecnolog\'{\i}as de \'ultima generaci\'on:
\begin{itemize}
\item \textbf{Go 1.24.3}: Alto rendimiento y concurrencia
\item \textbf{Rust 2021}: Seguridad de memoria y rendimiento
\item \textbf{Angular 16+}: Framework frontend moderno
\item \textbf{Tauri 2.0}: Aplicaciones de escritorio ligeras
\item \textbf{gRPC}: Comunicaci\'on entre servicios
\item \textbf{WebSockets}: Tiempo real
\end{itemize}

\newpage
